\chapter{Optimization Derivation: Expected Shortfall} \label{app:es_derivation}

In this appendix, we derive the optimization problems used in the main text,
building up from the Expected Shortfall objective to the implementable
quadratic programs solved in the empirical analysis.
In \Cref{app:es_auxiliary_formulation} we follow \cite{RockafellarUryasev2000} to derive the equivalence between direct
minimization of Expected Shortfall and an auxiliary formulation suitable for linear programming.
Following \cite{RegularizedMultiCurrencyExpectedShortfallPortfolios}, in \Cref{app:linear-program-representation}
we embed this equivalence in a mean-ES objective,
by replacing the conditional expectation with its empirical counterpart, and linearizing the
terms to assemble the base unregularized portfolio optimization problem as a linear
program.
In \Cref{app:regularization} we introduce three regularization techniques, an $L_1$
penalty, an $L_2$ penalty, and a leverage constraint, giving us the
base regularized quadratic program for portfolio optimization.
In \Cref{app:separate-optimization} we extend the framework to the multi-currency setting
under the separate (currency overlay) approach, deriving the asset and the currency quadratic programs.
Finally, in \Cref{app:joint-optimization} we derive the joint quadratic program formulation,
in which asset and currency weights are optimized simultaneously
through a single Expected Shortfall term that captures the combined return distribution.

\section{Derivation of the Auxiliary Expected Shortfall Formulation} \label{app:es_auxiliary_formulation}

We derive the equivalence between the direct minimization of Expected Shortfall and the auxiliary formulation proposed by Rockafellar and Uryasev \cite{RockafellarUryasev2000} that allows for formulating the portfolio optimization as a linear program.

We assume that the loss distribution is continuous, meaning the loss cumulative distribution function has no jumps and the probability density function $f(l)$ exists. Under these assumptions, Value-at-Risk (VaR) is uniquely defined \cite{QuantitativeRiskManagement2016}.

Define the loss random variable at time $t+1$, given portfolio weights $\bm{x}_t$, as $L(\bm{x}_t) = -\bm{x}_t^T \bm{R}_{t+1}$. Let $f(l)$ denote the probability density function of the loss $L(\bm{x}_t)$.\ The VaR at confidence level $\beta$ is defined as the $\beta$-quantile of the loss distribution
\begin{equation*}
    \text{VaR}_\beta(\bm{x}_t) = \min \{ \alpha : \mathbb{P}(L(\bm{x}_t) \le \alpha) \ge \beta \}.
\end{equation*}

The Expected Shortfall at confidence level $\beta$ is defined as the conditional expectation of losses exceeding $\text{VaR}_\beta$
\begin{equation*}
    \text{ES}_\beta(\bm{x}_t) = \mathbb{E}[L(\bm{x}_t) \mid L(\bm{x}_t) \ge \text{VaR}_\beta(\bm{x}_t)] = \frac{1}{1-\beta} \int_{\text{VaR}_\beta(\bm{x}_t)}^\infty l f(l) \, dl.
\end{equation*}

We now consider the auxiliary function $F_\beta$ defined as
\begin{equation*}
    F_\beta(\bm{x}_t, \alpha) = \alpha + \frac{1}{1-\beta} \mathbb{E}_t[(L(\bm{x}_t) - \alpha)^+].
\end{equation*}


We can rewrite the expectation term using the integral form
\begin{equation*}
    \mathbb{E}_t[(L(\bm{x}_t) - \alpha)^+] = \int_{-\infty}^\infty (l - \alpha)^+ f(l) \, dl = \int_{\alpha}^\infty (l - \alpha) f(l) \, dl.
\end{equation*}

Before minimizing $F_\beta(\bm{x}_t, \alpha)$ with respect to $\alpha$, we verify convexity. Using Leibniz's rule, the first and second derivative of $F_\beta$ with respect to $\alpha$ is
\begin{align*}
    \frac{\partial F_\beta}{\partial \alpha}
        &= 1 + \frac{1}{1-\beta} \frac{\partial}{\partial \alpha} \int_\alpha^\infty (l - \alpha) f(l) \, dl \\
    &= 1 + \frac{1}{1-\beta} \left[ -f(\alpha) \cdot 0 + \int_\alpha^\infty (-1) f(l) \, dl \right] \\
    &= 1 - \frac{1}{1-\beta} \int_{\alpha}^\infty f(l) \, dl \\
    \frac{\partial^2 F_\beta}{\partial \alpha^2} &= \frac{1}{1-\beta} f(\alpha).
\end{align*}
Since the probability density function $f(\alpha) \ge 0$ and $\beta < 1$, the second derivative is non-negative, $\frac{\partial^2 F_\beta}{\partial \alpha^2} \ge 0$. Thus, $F_\beta(\bm{x}_t, \alpha)$ is convex in $\alpha$, ensuring that the first-order condition is necessary and sufficient for a global minimum.

Differentiation of $F_\beta(\bm{x}_t, \alpha)$ with respect to $\alpha$ gives us the first-order condition for optimality
\begin{align*}
    \frac{\partial F_\beta}{\partial \alpha}
    &= 1 - \frac{1}{1-\beta} \int_\alpha^\infty f(l) \, dl \\
    &= 1 - \frac{1}{1-\beta} \mathbb{P}(L(\bm{x}_t) \ge \alpha).
\end{align*}

Setting the derivative to zero to find the optimal $\alpha^*$
\begin{align*}
    1 - \frac{1}{1-\beta} \mathbb{P}(L(\bm{x}_t) \ge \alpha^*) &= 0 \\
    \mathbb{P}(L(\bm{x}_t) \ge \alpha^*) &= 1 - \beta \\
    \mathbb{P}(L(\bm{x}_t) \le \alpha^*) &= \beta.
\end{align*}

This implies that the optimal $\alpha^*$ that minimizes the auxiliary function for a fixed $\bm{x}_t$ is exactly the Value-at-Risk at level $\beta$, so $\alpha^* = \text{VaR}_\beta(\bm{x}_t)$.

Substituting $\alpha^* = \text{VaR}_\beta(\bm{x}_t)$ back into $F_\beta(\bm{x}_t, \alpha)$ we obtain
\begin{align*}
    F_\beta(\bm{x}_t, \text{VaR}_\beta(\bm{x}_t)) &= \text{VaR}_\beta(\bm{x}_t) + \frac{1}{1-\beta} \int_{\text{VaR}_\beta(\bm{x}_t)}^\infty (l - \text{VaR}_\beta(\bm{x}_t)) f(l) \, dl \\
    &= \text{VaR}_\beta(\bm{x}_t) + \frac{1}{1-\beta} \left( \int_{\text{VaR}_\beta(\bm{x}_t)}^\infty l f(l) \, dl - \text{VaR}_\beta(\bm{x}_t) \int_{\text{VaR}_\beta(\bm{x}_t)}^\infty f(l) \, dl \right).
\end{align*}

Recognizing that $\int_{\text{VaR}_\beta(\bm{x}_t)}^\infty f(l) \, dl = \mathbb{P}(L(\bm{x}_t) \ge \text{VaR}_\beta(\bm{x}_t)) = 1-\beta$, we have
\begin{align*}
    F_\beta(\bm{x}_t, \text{VaR}_\beta(\bm{x}_t)) &= \text{VaR}_\beta(\bm{x}_t) + \frac{1}{1-\beta} \int_{\text{VaR}_\beta(\bm{x}_t)}^\infty l f(l) \, dl - \frac{1}{1-\beta} \text{VaR}_\beta(\bm{x}_t) (1-\beta) \\
    &= \text{VaR}_\beta(\bm{x}_t) + \text{ES}_\beta(\bm{x}_t) - \text{VaR}_\beta(\bm{x}_t) \\
    &= \text{ES}_\beta(\bm{x}_t).
\end{align*}

Thus, minimizing $\text{ES}_\beta(\bm{x}_t)$ with respect to $\bm{x}_t$ is equivalent to minimizing $F_\beta(\bm{x}_t, \alpha)$ jointly with respect to $\bm{x}_t$ and $\alpha$
\begin{equation} \label{eq:joint_opt}
    \min_{\bm{x}_t} \text{ES}_\beta(\bm{x}_t) = \min_{\bm{x}_t, \alpha} \left( \alpha + \frac{1}{1-\beta} \mathbb{E}_t[(-\bm{x}_t^T \bm{R}_{t+1} - \alpha)^+] \right).
\end{equation}

Finally, we observe that $F_\beta(\bm{x}_t, \alpha)$ is jointly convex in $(\bm{x}_t, \alpha)$. The term $\alpha$ is linear (hence convex). The inner map $(\bm{x}_t, \alpha) \mapsto (-\bm{x}_t^\top \bm{R}_{t+1} - \alpha)$ is affine. The outer function $z \mapsto z^+$ is convex and nondecreasing. The composition of a convex function with an affine map is convex \cite{boyd2004convex}. Thus $(-\bm{x}_t^\top \bm{R}_{t+1} - \alpha)^+$ is jointly convex in $(\bm{x}_t, \alpha)$. Taking the expectation preserves convexity.


Since the objective function is convex, any local minimum is global, ensuring the optimization problem can be solved reliably with standard convex optimization methods.

\section{Linear Program Representation of Portfolio Optimization} \label{app:linear-program-representation} % another possible title: Linear program representation of portfolio optimization

The equivalence in \eqref{eq:joint_opt} reduces Expected Shortfall minimization to the
minimization of the jointly convex auxiliary function $F_\beta(\bm{x}_t, \alpha)$
\begin{equation*}
    \min_{\bm{x}_t} \text{ES}_\beta(\bm{x}_t)
    = \min_{\bm{x}_t, \alpha} \left( \alpha + \frac{1}{1-\beta}
    \mathbb{E}_t[(-\bm{x}_t^\top \bm{R}_{t+1} - \alpha)^+] \right).
\end{equation*}
This formulation is not yet appropriate for portfolio optimization, as it only optimizes for downside risk, without
taking into account the expected return of the portfolio. It also depends on the unknown
expectation $\mathbb{E}_t[\cdot]$. We address this by incorporating the expected return into the optimization
objective and replacing the expectation with its empirical estimator.

Following \cite{goldberg2013minimizingshortfall} and \cite{RegularizedMultiCurrencyExpectedShortfallPortfolios}, we trade off expected return against %
downside risk through a relative risk-aversion coefficient $\gamma > 0$
\begin{equation} \label{eq:mean_es}
    \bm{x}_t^* = \argmin_{\bm{x}_t} \left\{ -\bm{\mu}_t^\top \bm{x}_t + \frac{\gamma}{2}
    \text{ES}_\beta(\bm{x}_t) \right\}, \quad \text{s.t. } \bm{x}_t^\top \bm{1}_N = 1,
\end{equation}
where $\bm{\mu}_t$ denotes the asset return mean estimator and the budget constraint
$\bm{x}_t^\top \bm{1}_N = 1$ enforces full investment. Substituting the auxiliary
representation \eqref{eq:joint_opt} into \eqref{eq:mean_es} gives
\begin{equation*}
    \bm{x}_t^*, \alpha^* = \argmin_{\bm{x}_t, \alpha} \left\{ -\bm{\mu}_t^\top \bm{x}_t
    + \frac{\gamma}{2} \alpha
    + \frac{\gamma}{2(1-\beta)} \mathbb{E}_t[(-\bm{x}_t^\top \bm{R}_{t+1} - \alpha)^+] \right\}.
\end{equation*}

Let $\bm{y}_1, \dots, \bm{y}_q \in \mathbb{R}^N$ denote $q$ historical realizations of the joint
return vector $\bm{R}_{t+1}$.\ We approximate the conditional expectation by the empirical
average over this sample
\begin{equation*}
    \mathbb{E}_t[(-\bm{x}_t^\top \bm{R}_{t+1} - \alpha)^+]
    \approx \frac{1}{q} \sum_{k=1}^q (-\bm{x}_t^\top \bm{y}_k - \alpha)^+,
\end{equation*}
so that the portfolio optimization problem becomes
\begin{equation} \label{eq:empirical_es}
    \bm{x}_t^*, \alpha^* = \argmin_{\bm{x}_t, \alpha} \left\{ -\bm{\mu}_t^\top \bm{x}_t
    + \frac{\gamma}{2} \alpha
    + \frac{\gamma}{2q(1-\beta)} \sum_{k=1}^q (-\bm{x}_t^\top \bm{y}_k - \alpha)^+ \right\},
    \quad \text{s.t. } \bm{x}_t^\top \bm{1}_N = 1.
\end{equation}

The summand $(-\bm{x}_t^\top \bm{y}_k - \alpha)^+$ is convex but non-linear. We linearize it by
introducing, for each observation $k$, an auxiliary variable $u_k \ge 0$ subject to
\begin{equation} \label{eq:slack_constraint}
    u_k \ge -\bm{x}_t^\top \bm{y}_k - \alpha
    \qquad \Longleftrightarrow \qquad
    \bm{x}_t^\top \bm{y}_k + \alpha + u_k \ge 0.
\end{equation}
Since the objective in \eqref{eq:empirical_es} is strictly increasing in each $u_k$
(its coefficient $\tfrac{\gamma}{2q(1-\beta)} > 0$) and minimized, every $u_k$ is driven
to its lower bound at the optimum, so $u_k^* = (-\bm{x}_t^\top \bm{y}_k - \alpha)^+$. The slack
constraints therefore reproduce the term $(-\bm{x}_t^\top \bm{y}_k - \alpha)^+$ while keeping all
relations linear.


By collecting the decision variables in the vector
\begin{equation*}
    \tilde{\bm{x}}_t = \begin{bmatrix} x_{1,t}, \dots, x_{N,t}, &  \alpha, & u_1, \dots, u_q \end{bmatrix}^\top
    = \begin{bmatrix} \bm{x}_t^\top, & \alpha, & \bm{u}^\top \end{bmatrix}^\top \in \mathbb{R}^{(N+q+1)},
\end{equation*}
and constants in the vector
\begin{equation*}
    \bm{c} = \begin{bmatrix} -\bm{\mu}_t^\top, &  \dfrac{\gamma}{2}, & \dfrac{\gamma}{2q(1-\beta)} \bm{1}_q^\top \end{bmatrix}^\top \in \mathbb{R}^{(N+q+1)},
\end{equation*}
the portfolio optimization problem \eqref{eq:empirical_es} can be written as a linear program
\begin{equation} \label{eq:base_lp}
\begin{aligned}
    & \ \ \tilde{\bm{x}}_t^* = \argmin_{\tilde{\bm{x}}_t} \quad
     \bm{c}^\top \tilde{\bm{x}}_t \\[4pt]
    \text{s.t.} \quad
    & \bm{D} \tilde{\bm{x}}_t =
    \begin{bmatrix} \bm{1}_N^\top & \bm{0}_{q+1}^\top \end{bmatrix}
    \tilde{\bm{x}}_t = E = 1, \\[2pt]
    & \bm{A} \tilde{\bm{x}}_t =
    \begin{bmatrix} -\bm{R}^{\text{data}} & -\bm{1}_q & -\bm{I}_q \end{bmatrix}
    \tilde{\bm{x}}_t \le \bm{b} = \bm{0}_q, \\[2pt]
    & \alpha \ge 0, \quad u_k \ge 0 \quad \text{for } k = 1, \dots, q.
\end{aligned}
\end{equation}
Here $\bm{D} \in \mathbb{R}^{1 \times (N+q+1)}$ encodes the budget constraint, while
$\bm{A} \in \mathbb{R}^{q \times (N+q+1)}$ and $\bm{b} \in \mathbb{R}^{q}$ encode the
$q$ inequality constraints arising from the linearization \eqref{eq:slack_constraint}. The matrix $\bm{R}^{\text{data}} \in \mathbb{R}^{q \times N}$ is comprised of
sampled return vectors $\bm{y}_k^\top$ row by row, $\bm{0}_q$ is the $q$-dimensional zero vector,
and $\bm{I}_q$ the $q$-dimensional identity matrix.

Problem \eqref{eq:base_lp} has a linear objective subject to one linear equality and $q$
linear inequalities, and is thus a linear program solvable with standard methods. This
formulation serves as the base unregularized portfolio optimization problem under Expected Shortfall.
Regularization terms and additional constraints are added in \Cref{app:regularization} of the appendix.
\section{Regularization} \label{app:regularization}

The base problem \eqref{eq:base_lp} is known to be numerically unstable in
high-dimensional settings, producing extreme and concentrated weights \citep{ciliberti2007feasibility, kondor2007noise}. We stabilize it
with three regularization techniques. An $L_1$-norm penalty that promotes sparsity, an $L_2$-norm penalty
that promotes shrinkage, and a leverage constraint that bounds gross exposure. Each is
appended to \eqref{eq:base_lp} while preserving its linear or quadratic structure. The
$L_1$ penalty and the leverage constraint share a common variable decomposition, which
we introduce first.

\subsection{$L_1$ Regularization} \label{app:l1}

The $L_1$ penalty $\lambda_1 \|\bm{x}_t\|_1 = \lambda_1 \sum_{i=1}^N |x_{i,t}|$, with
$\lambda_1 \ge 0$, is non-linear through the absolute value. Following \cite{roncalli2013introduction},
we linearize it by splitting each weight into its positive and negative parts via the
auxiliary variables $\Delta_{i,t}^+ \ge 0$ and $\Delta_{i,t}^- \ge 0$
\begin{equation} \label{eq:l1_decomp}
    x_{i,t} = \Delta_{i,t}^+ - \Delta_{i,t}^-,
    \qquad
    |x_{i,t}| = \Delta_{i,t}^+ + \Delta_{i,t}^-.
\end{equation}
The second identity holds at the optimum. Minimizing the $L_1$ penalty drives the parts to
their smallest feasible values, so at most one of $\Delta_{i,t}^+$, $\Delta_{i,t}^-$ is
non-zero, giving $\Delta_{i,t}^+ = \max(x_{i,t}, 0)$ and
$\Delta_{i,t}^- = \max(-x_{i,t}, 0)$. The $L_1$ penalty is then linear in the auxiliary
variables
\begin{equation*}
    \lambda_1 \|\bm{x}_t\|_1 = \lambda_1 \bm{1}_N^\top (\bm{\Delta}_t^+ + \bm{\Delta}_t^-).
\end{equation*}

We extend the decision vector to carry the two parts
\begin{equation*}
    \tilde{\bm{x}}_t = \begin{bmatrix} \bm{x}_t^\top,& \alpha,& \bm{u}^\top,& (\bm{\Delta}_t^-)^\top,& (\bm{\Delta}_t^+)^\top \end{bmatrix}^\top
    \in \mathbb{R}^{3N+q+1},
\end{equation*}
and append the penalty coefficients to the cost vector
\begin{equation*}
    \tilde{\bm{c}} = \begin{bmatrix} -\bm{\mu}_t^\top,& \tfrac{\gamma}{2},& \tfrac{\gamma}{2q(1-\beta)} \bm{1}_q^\top,& \lambda_1 \bm{1}_N^\top,& \lambda_1 \bm{1}_N^\top \end{bmatrix}^\top.
\end{equation*}
The decomposition \eqref{eq:l1_decomp} is enforced as $N$ additional linear equality
constraints, $\bm{x}_t + \bm{\Delta}_t^- - \bm{\Delta}_t^+ = \bm{0}_N$, alongside the original budget
constraint. All relations remain linear, so the optimization problem is still a linear program after the addition of the $L_1$ penalty.

\subsection{$L_2$ Regularization} \label{app:l2}

The $L_2$ penalty $\lambda_2 \|\bm{x}_t\|_2^2 = \lambda_2 \bm{x}_t^\top \bm{x}_t$, with
$\lambda_2 \ge 0$, is quadratic in the weights and is written as a quadratic form
\begin{equation*}
    \lambda_2 \|\bm{x}_t\|_2^2 = \tfrac{1}{2}\, \bm{x}_t^\top \bm{Q}\, \bm{x}_t,
    \qquad \bm{Q} = 2 \lambda_2 \bm{I}_N.
\end{equation*}
Embedding this in the augmented vector $\tilde{\bm{x}}_t$, the penalty becomes
$\tfrac{1}{2}\, \tilde{\bm{x}}_t^\top \tilde{\bm{Q}}\, \tilde{\bm{x}}_t$ with $\tilde{\bm{Q}}$ placing the
block $2\lambda_2 \bm{I}_N$ on the asset weights and zeros on the auxiliary variables
\begin{equation*}
    \tilde{\bm{Q}} =
    \begin{bmatrix}
        2 \lambda_2 \bm{I}_N & \bm{0}_{N \times (2N+q+1)} \\
        \bm{0}_{(2N+q+1) \times N} & \bm{0}_{(2N+q+1) \times (2N+q+1)}
    \end{bmatrix}.
\end{equation*}
Adding this term turns the $L_1-$augmented linear objective of the problem \eqref{eq:base_lp} into a quadratic objective
$\tilde{\bm{c}}^\top \tilde{\bm{x}}_t + \tfrac{1}{2}\, \tilde{\bm{x}}_t^\top \tilde{\bm{Q}}\, \tilde{\bm{x}}_t$.
Since $\tilde{\bm{Q}}$ is positive semi-definite, the objective remains convex, and the optimization problem becomes a
quadratic program with linear constraints.

\subsection{Leverage Constraint} \label{app:leverage}

The leverage constraint $\|\bm{x}_t\|_1 \le \ell$, with $\ell \ge 1$, bounds the portfolio's
gross exposure and guarantees a bounded feasible region even when $\lambda_1 = 0$. Reusing
the decomposition \eqref{eq:l1_decomp}, it is expressed as a single linear inequality
\begin{equation} \label{eq:leverage_constraint}
    \|\bm{x}_t\|_1 = \bm{1}_N^\top (\bm{\Delta}_t^+ + \bm{\Delta}_t^-) \le \ell.
\end{equation}
The choice $\ell = 1$ rules out leverage entirely and is equivalent to a hard $L_1$
constraint. As a linear inequality, it preserves the quadratic program structure.

\subsection{Quadratic Program Representation of the Regularized Portfolio Optimization} \label{app:reg_qp}

Adding all three regularization terms to the base problem \eqref{eq:base_lp}, makes the portfolio optimization problem under Expected Shortfall
a quadratic program with linear constraints
\begin{equation} \label{eq:reg_qp}
\begin{aligned}
    \tilde{\bm{x}}_t^* &= \argmin_{\tilde{\bm{x}}_t} \quad
    \tilde{\bm{c}}^\top \tilde{\bm{x}}_t + \tfrac{1}{2}\, \tilde{\bm{x}}_t^\top \tilde{\bm{Q}}\, \tilde{\bm{x}}_t, \\[2pt]
    \text{s.t.} \quad
     \tilde{\bm{D}} \tilde{\bm{x}}_t &=
    \begin{bmatrix} \bm{1}_N^\top & \bm{0}_{q+1}^\top & \bm{0}_N^\top & \bm{0}_N^\top \\
    \bm{I}_N & \bm{0}_{N \times (q+1)} & \bm{I}_N & -\bm{I}_N \end{bmatrix} \tilde{\bm{x}}_t
    = \tilde{\bm{E}} = \begin{bmatrix} 1 \\ \bm{0}_N \end{bmatrix}, \\[6pt]
     \tilde{\bm{A}} \tilde{\bm{x}}_t &=
    \begin{bmatrix} -\bm{R}^{\text{data}} & -\bm{1}_q & -\bm{I}_q & \bm{0}_{q \times N} & \bm{0}_{q \times N} \\
    \bm{0}_N^\top & 0 & \bm{0}_q^\top & \bm{1}_N^\top & \bm{1}_N^\top \end{bmatrix} \tilde{\bm{x}}_t
    \le \bm{b} = \begin{bmatrix} \bm{0}_q \\ \ell \end{bmatrix}, \\[6pt]
     \alpha \ge 0, & \quad u_k \ge 0,\quad \Delta_{i,t}^- \ge 0,\quad \Delta_{i,t}^+ \ge 0
    \quad \text{for } i = 1, \dots, N,\ k = 1, \dots, q.
\end{aligned}
\end{equation}
The first row of $\tilde{\bm{D}}$ enforces the budget constraint and the remaining $N$ rows
enforce the weight decomposition \eqref{eq:l1_decomp}. The first $q$ rows of $\tilde{\bm{A}}$
encode the inequality constraints from the linearization \eqref{eq:slack_constraint}, and are extended with zeros for the new
variables, and the final row of $\tilde{\bm{A}}$ imposes the leverage constraint \eqref{eq:leverage_constraint}.

\section{Separate (Currency Overlay) Optimization} \label{app:separate-optimization}
In the separate (currency overlay) portfolio optimization approach, the asset and currency allocations are
optimized in two sequential steps. First, the asset weights $\bm{x}_t \in \mathbb{R}^N$ are
optimized on the fully hedged returns, then, conditional on the optimal $\bm{x}_t^*$, the
currency exposures $\bm{\psi}_t \in \mathbb{R}^M$ are optimized.
Let
\begin{align*}
\bm{\mu}^{fh} &= \mathbb{E}_t[\bm{R}^{fh}_{t+1}] \in \mathbb{R}^N, \\
\bm{\mu}^c &= \mathbb{E}_t[\bm{e}_{t+1} - \bm{f}_t] \in \mathbb{R}^M,
\end{align*}
denote the expected fully hedged asset return and the expected
excess currency return respectively. The two otimization steps are then given by
\begin{equation} \label{eq:sep_asset_obj}
\begin{gathered}
    \bm{x}_t^* = \argmin_{\bm{x}_t} \left\{ -\bm{x}_t^\top \bm{\mu}^{fh}
    + \tfrac{\gamma}{2} \text{ES}_\beta\left(\bm{x}_t^\top \bm{R}^{fh}_{t+1}\right)
    + \lambda_1^a \|\bm{x}_t\|_1 + \lambda_2^a \|\bm{x}_t\|_2^2 \right\} \\[4pt]
    \text{s.t. } \bm{x}_t^\top \bm{1}_N = 1,\ \|\bm{x}_t\|_1 \le \ell,
\end{gathered}
\end{equation}
\begin{equation} \label{eq:sep_curr_obj}
\begin{split}
    \bm{\psi}_t^* \mid \bm{x}_t^* = \argmin_{\bm{\psi}_t} \Big\{\, &
    {-}\bm{\psi}_t^\top \bm{\mu}^c
    + \tfrac{\gamma}{2}\, \text{ES}_\beta\!\left( \bm{x}_t^{*\top} \bm{R}^{fh}_{t+1}
    + \bm{\psi}_t^\top (\bm{e}_{t+1} - \bm{f}_t) \right) \\
    & + \lambda_1^c \|\bm{\psi}_t\|_1 + \lambda_2^c \|\bm{\psi}_t\|_2^2 \,\Big\}.
\end{split}
\end{equation}
The regularization parameters are kept separate across assets ($\lambda_1^a, \lambda_2^a$)
and currencies ($\lambda_1^c, \lambda_2^c$).

\subsection{Asset Optimization} \label{app:sep_asset}

The asset optimization step \eqref{eq:sep_asset_obj} is identical to the regularized
single-currency program \eqref{eq:reg_qp}. It carries a budget constraint, $L_1$
and $L_2$ penalty terms, and a leverage constraint.
It is therefore obtained directly from \eqref{eq:reg_qp} by substituting the matrix of historically sampled fully hedged
returns $\bm{R}^{fh}\in \mathbb{R}^{q \times N}$ for $\bm{R}^{\text{data}}$, the mean $\bm{\mu}^{fh}$ for $\bm{\mu}_t$, and the
asset-specific penalties $\lambda_1^a, \lambda_2^a$. With the augmented vector
\begin{equation*}
    \tilde{\bm{x}}_t = \begin{bmatrix} \bm{x}_t^\top,& \alpha_x,& \bm{u}_x^\top,&
    (\bm{\Delta}_{x,t}^-)^\top,& (\bm{\Delta}_{x,t}^+)^\top \end{bmatrix}^\top \in \mathbb{R}^{3N+q+1},
\end{equation*}
the cost vector and penalty matrix are
\begin{equation*}
    \tilde{\bm{c}} =
    \begin{bmatrix}
        -(\bm{\mu}^{fh})^\top, & \tfrac{\gamma}{2}, &
        \tfrac{\gamma}{2q(1-\beta)} \bm{1}_q^\top, &
        \lambda_1^a \bm{1}_N^\top, & \lambda_1^a \bm{1}_N^\top
    \end{bmatrix}^\top,
\end{equation*}

\begin{equation*}
    \tilde{\bm{Q}} =
    \begin{bmatrix}
        2 \lambda_2^a \bm{I}_N & \bm{0}_{N \times (2N+q+1)} \\
        \bm{0}_{(2N+q+1) \times N} & \bm{0}_{(2N+q+1) \times (2N+q+1)}
    \end{bmatrix},
\end{equation*}
and the equality, inequality, and sign constraints are exactly those of
\eqref{eq:reg_qp} with $\bm{R}^{fh}$ in place of $\bm{R}^{\text{data}}$. The step yields the optimal
asset weights $\bm{x}_t^*$, which are then fixed in the currency overlay optimization.

\subsection{Currency Overlay Optimization} \label{app:sep_currency}

Conditional on $\bm{x}_t^*$, the currency step \eqref{eq:sep_curr_obj} minimizes the Expected
Shortfall of the combined return $\bm{x}_t^{*\top} \bm{R}^{fh}_{t+1} + \bm{\psi}_t^\top (\bm{e}_{t+1} - \bm{f}_t)$
over the currency exposures $\bm{\psi}_t$. Applying the same auxiliary representation and the empirical approximation as in
\Cref{app:linear-program-representation}, we introduce the auxiliary variable
$\alpha_\psi$ and slack variables $u_{k,\psi} \ge 0$. Since $\bm{x}_t^*$ is now fixed, the
slack constraint becomes
\begin{equation*}
    u_{k,\psi} \ge -\bm{\psi}_t^\top (\bm{e}_k - \bm{f}_t) - \bm{x}_t^{*\top} \bm{R}^{fh}_k - \alpha_\psi
    \quad \Longleftrightarrow \quad
    -\bm{\psi}_t^\top (\bm{e}_k - \bm{f}_t) - \alpha_\psi - u_{k,\psi} \le \bm{x}_t^{*\top} \bm{R}^{fh}_k,
\end{equation*}
where $\bm{R}^{fh}_k$ is the $k$-th sampled fully hedged return vector and $\bm{e}_k$ the
$k$-th sampled currency return vector. The fixed asset-portfolio
returns therefore appear on the right-hand side. We stack the sampled excess currency returns
$(\bm{e}_k - \bm{f}_t)^\top$ row by row into $\bm{R}^c \in \mathbb{R}^{q \times M}$, and collect the
fixed asset-portfolio returns in $\bm{R}^{fh}_P = \bm{R}^{fh} \bm{x}_t^* \in \mathbb{R}^q$. We linearize the currency $L_1$ penalty with
the decomposition $\psi_{c,t} = \Delta_{\psi,c,t}^+ - \Delta_{\psi,c,t}^-$ as in
\Cref{app:l1}, and form the augmented vector
\begin{equation*}
    \tilde{\bm{\psi}}_t = \begin{bmatrix} \bm{\psi}_t^\top,& \alpha_\psi,& \bm{u}_\psi^\top,&
    (\bm{\Delta}_{\psi,t}^-)^\top,& (\bm{\Delta}_{\psi,t}^+)^\top \end{bmatrix}^\top \in \mathbb{R}^{3M+q+1},
\end{equation*}
The cost vector and penalty matrix are
\begin{equation*}
    \tilde{\bm{c}} = \begin{bmatrix} -(\bm{\mu}^c)^\top,& \tfrac{\gamma}{2},& \tfrac{\gamma}{2q(1-\beta)} \bm{1}_q^\top,& \lambda_1^c \bm{1}_M^\top,& \lambda_1^c \bm{1}_M^\top \end{bmatrix}^\top,
\end{equation*}
\begin{equation*}
    \tilde{\bm{Q}} =
    \begin{bmatrix}
        2 \lambda_2^c \bm{I}_M & \bm{0}_{M \times (2M+q+1)} \\
        \bm{0}_{(2M+q+1) \times M} & \bm{0}_{(2M+q+1) \times (2M+q+1)}
    \end{bmatrix}.
\end{equation*}
The currency overlay optimization is then the quadratic program
\begin{equation} \label{eq:sep_curr_qp}
\begin{aligned}
    \bm{\psi}_t^* &= \argmin_{\tilde{\bm{\psi}}_t} \quad
     \tilde{\bm{c}}^\top \tilde{\bm{\psi}}_t + \tfrac{1}{2}\, \tilde{\bm{\psi}}_t^\top \tilde{\bm{Q}}\, \tilde{\bm{\psi}}_t, \\[2pt]
    \text{s.t.} \quad
    \tilde{\bm{D}} \tilde{\bm{\psi}}_t &=
    \begin{bmatrix} \bm{I}_M & \bm{0}_{M \times (q+1)} & \bm{I}_M & -\bm{I}_M \end{bmatrix} \tilde{\bm{\psi}}_t
    = \tilde{\bm{E}} = \bm{0}_M, \\[4pt]
    \tilde{\bm{A}} \tilde{\bm{\psi}}_t &=
    \begin{bmatrix} -\bm{R}^c & -\bm{1}_q & -\bm{I}_q & \bm{0}_{q \times M} & \bm{0}_{q \times M} \end{bmatrix}
    \tilde{\bm{\psi}}_t \le \bm{b} = \bm{R}^{fh}_P, \\[4pt]
    \alpha_\psi &\ge 0, \quad u_{k,\psi} \ge 0, \quad
    \Delta_{\psi,c,t}^- \ge 0, \quad \Delta_{\psi,c,t}^+ \ge 0
    \quad \text{for } c = 1, \dots, M,\ k = 1, \dots, q.
\end{aligned}
\end{equation}

The currency optimization differs from the asset optimization \eqref{eq:sep_asset_obj} in the following ways.
First, there is no budget constraint, the currency portfolio is a
zero-investment portfolio, so the domestic exposure $\psi_{0,t}$ is inferred residually and
the equality block $\tilde{\bm{D}}$ enforces only the weight decomposition, with right-hand
side $\tilde{\bm{E}} = \bm{0}_M$ rather than $[1, \bm{0}_M^\top]^\top$. Second, the fixed asset leg
shifts the inequality right-hand side from $\bm{0}_q$ to $\bm{R}^{fh}_P = \bm{R}^{fh} \bm{x}_t^*$. Third, no
leverage constraint is imposed on the currency exposures through matrix $\tilde{\bm{A}}$. The $L_1$ and $L_2$ regularizing terms are
as before, through the last two blocks of $\tilde{\bm{c}}$ and the main block of
$\tilde{\bm{Q}}$, so the currency overlay remains a convex quadratic program.

\section{Joint Optimization} \label{app:joint-optimization}

In the joint portfolio optimization approach, asset and currency weights are optimized
simultaneously rather than sequentially. We stack them into the combined decision vector
$\bm{\theta}_t = [\bm{x}_t^\top, \bm{\psi}_t^\top]^\top \in \mathbb{R}^{N+M}$ and the combined return vector
$\bm{R}_{t+1} = [\bm{R}^{fh\top}_{t+1},\, (\bm{e}_{t+1} - \bm{f}_t)^\top]^\top \in \mathbb{R}^{N+M}$, with joint mean
$\bm{\mu} = \mathbb{E}_t[\bm{R}_{t+1}] \in \mathbb{R}^{N+M}$. The joint optimization problem is
\begin{equation} \label{eq:joint_obj}
\begin{gathered}
    \bm{\theta}_t^* = \argmin_{\bm{\theta}_t} \left\{ -\bm{\theta}_t^\top \bm{\mu}
    + \tfrac{\gamma}{2} \text{ES}_\beta\left(\bm{\theta}_t^\top \bm{R}_{t+1}\right)
    + \|\bm{S}_1 \bm{\theta}_t\|_1 + \|\bm{S}_2 \bm{\theta}_t\|_2^2 \right\}, \\[4pt]
    \text{s.t. } \bm{\theta}_t^\top \bm{q}_{N,M} = 1, \quad \|\bm{L} \bm{\theta}_t\|_1 \le \ell.
\end{gathered}
\end{equation}
The following matrices and vector adapt the problem to the
joint optimization setting
\begin{equation*}
    \bm{q}_{N,M} = \begin{bmatrix} \bm{1}_N \\ \bm{0}_M \end{bmatrix}, \quad
    \bm{L} = \begin{bmatrix} \bm{I}_N & \bm{0}_{N \times M} \\ \bm{0}_{M \times N} & \bm{0}_{M \times M} \end{bmatrix}, \quad
    \bm{S}_1 = \begin{bmatrix} \lambda_1^a \bm{I}_N & 0 \\ 0 & \lambda_1^c \bm{I}_M \end{bmatrix}, \quad
    \bm{S}_2 = \begin{bmatrix} \sqrt{\lambda_2^a}\, \bm{I}_N & 0 \\ 0 & \sqrt{\lambda_2^c}\, \bm{I}_M \end{bmatrix}.
\end{equation*}
The vector $\bm{q}_{N,M}$ applies the budget constraint to the asset weights only, so that
$\bm{\theta}_t^\top \bm{q}_{N,M} = \bm{x}_t^\top \bm{1}_N = 1$, leaving the currency weights as a
zero-investment overlay. The matrix $\bm{L}$ applies the leverage constraint
to the asset weights only, giving $\|\bm{L}\bm{\theta}_t\|_1 = \|\bm{x}_t\|_1$. The diagonal regularization
matrices $\bm{S}_1$ and $\bm{S}_2$ allow the asset and currency penalties to differ within a
single objective,
\begin{equation*}
    \|\bm{S}_1 \bm{\theta}_t\|_1 = \lambda_1^a \|\bm{x}_t\|_1 + \lambda_1^c \|\bm{\psi}_t\|_1,
    \qquad
    \|\bm{S}_2 \bm{\theta}_t\|_2^2 = \lambda_2^a \|\bm{x}_t\|_2^2 + \lambda_2^c \|\bm{\psi}_t\|_2^2.
\end{equation*}

Applying the auxiliary representation of \Cref{app:es_auxiliary_formulation} and the
empirical approximation of \Cref{app:linear-program-representation} to the joint return
$\bm{\theta}_t^\top \bm{R}_{t+1}$, we introduce the auxiliary variable $\alpha$ and slack
variables $u_k \ge 0$ for $k = 1, \dots, q$, and linearize the $L_1$ penalty by
decomposing all $N+M$ components of $\bm{\theta}_t$ into positive and negative parts as in
\Cref{app:l1}. Collecting the asset and currency parts separately, we form the
augmented vector
\begin{equation*}
    \tilde{\bm{\theta}}_t = \begin{bmatrix}\, \bm{x}_t^\top,& \bm{\psi}_t^\top,& \alpha,& \bm{u}^\top,&
    (\bm{\Delta}_{x,t}^-)^\top,& (\bm{\Delta}_{\psi,t}^-)^\top,&
    (\bm{\Delta}_{x,t}^+)^\top,& (\bm{\Delta}_{\psi,t}^+)^\top \,\end{bmatrix}^\top
    \in \mathbb{R}^{3N+3M+q+1}.
\end{equation*}
Let $\bm{R}^{\text{all}} = [\, \bm{R}^{fh}\ \ \bm{R}^c \,] \in \mathbb{R}^{q \times (N+M)}$ denote the
matrix of jointly sampled asset and currency returns. The cost vector and penalty matrix are
\begin{equation*}
    \tilde{\bm{c}} = \begin{bmatrix}\, -\bm{\mu}^\top,& \tfrac{\gamma}{2},&
    \tfrac{\gamma}{2q(1-\beta)} \bm{1}_q^\top,&
    \lambda_1^a \bm{1}_N^\top,& \lambda_1^c \bm{1}_M^\top,&
    \lambda_1^a \bm{1}_N^\top,& \lambda_1^c \bm{1}_M^\top \,\end{bmatrix}^\top,
\end{equation*}
\begin{equation*}
    \tilde{\bm{Q}} =
    \begin{bmatrix}
        2 \lambda_2^a \bm{I}_N & 0 & 0 \\
        0 & 2 \lambda_2^c \bm{I}_M & 0 \\
        0 & 0 & \bm{0}_{(2N+2M+q+1) \times (2N+2M+q+1)}
    \end{bmatrix}.
\end{equation*}
The joint portfolio optimization problem is then the quadratic program
\begin{equation} \label{eq:joint_qp}
\begin{aligned}
    \tilde{\bm{\theta}}_t^* &= \argmin_{\tilde{\bm{\theta}}_t} \quad
    \tilde{\bm{c}}^\top \tilde{\bm{\theta}}_t + \tfrac{1}{2}\, \tilde{\bm{\theta}}_t^\top \tilde{\bm{Q}}\, \tilde{\bm{\theta}}_t, \\[2pt]
    \text{s.t.} \quad
    \tilde{\bm{D}} \tilde{\bm{\theta}}_t &=
    \begin{bmatrix} \bm{q}_{N,M}^\top & \bm{0}_{q+1}^\top & \bm{0}_{N+M}^\top & \bm{0}_{N+M}^\top \\
    \bm{I}_{N+M} & \bm{0}_{(N+M) \times (q+1)} & \bm{I}_{N+M} & -\bm{I}_{N+M} \end{bmatrix} \tilde{\bm{\theta}}_t
    = \tilde{\bm{E}} = \begin{bmatrix} 1 \\ \bm{0}_{N+M} \end{bmatrix}, \\[6pt]
    \tilde{\bm{A}} \tilde{\bm{\theta}}_t &=
    \begin{bmatrix} -\bm{R}^{\text{all}} & -\bm{1}_q & -\bm{I}_q & \bm{0}_{q \times 2N} & \bm{0}_{q \times 2M} \\
    \bm{0}_{N+M}^\top & \bm{0}_{q+1}^\top & \bm{q}_{N,M}^\top & & \bm{q}_{N,M}^\top \end{bmatrix} \tilde{\bm{\theta}}_t
    \le \bm{b} = \begin{bmatrix} \bm{0}_q \\ \ell \end{bmatrix}, \\[6pt]
    \alpha &\ge 0, \quad u_k \ge 0, \quad \text{for} \ k = 1, \dots, q, \\[2pt]
    \Delta_{x,i,t}^\pm &\ge 0, \quad \Delta_{\psi,c,t}^\pm \ge 0,
    \quad \text{for}\ i = 1, \dots, N, \ c = 1, \dots, M.
\end{aligned}
\end{equation}
In $\tilde{\bm{D}}$, the first row enforces the asset budget constraint $\bm{x}_t^\top \bm{1}_N = 1$,
and the remaining $N+M$ rows enforce the weight decomposition \eqref{eq:l1_decomp} for
all components of $\bm{\theta}_t$. In $\tilde{\bm{A}}$, the first $q$ rows encode the linearization
inequalities \eqref{eq:slack_constraint} applied to the joint return matrix $\bm{R}^{\text{all}}$.
The final row imposes the asset-only leverage constraint
\eqref{eq:leverage_constraint}, giving $\bm{1}_N^\top(\bm{\Delta}_{x,t}^+ + \bm{\Delta}_{x,t}^-) = \|\bm{x}_t\|_1 \le \ell$.

The joint formulation solves a single quadratic program over $\tilde{\bm{\theta}}_t$,
capturing the interaction between asset and currency exposures through the
joint Expected Shortfall term and the off-diagonal return covariation in $\bm{R}^{\text{all}}$.
This comes at the cost of a higher-dimensional problem
relative to the two smaller programs of the separate approach.

All quadratic programs derived in this appendix share the same convex structure and
are solvable with the same quadratic programming solver. The auxiliary representation of Expected Shortfall ensures
this uniformity. Each extension of the program, from the base regularized single currency
formulation \eqref{eq:reg_qp}, to the joint multi-currency program above \eqref{eq:joint_qp},
modifies only the return matrix and the constraint blocks that encode the budget,
leverage, and weight decomposition, without changing the program structure.
